\subsection{Calibración del costo de credenciales}\label{sec:calibracion-credenciales}

El 19 de septiembre de 2026 se repitió la medición del costo de un hash nuevo en un equipo AMD Ryzen 7 7730U, el mismo modelo de la medición histórica de la \cref{sec:pruebas-argon2}, con 16 CPU lógicas y Linux 7.1.13. El binario de desarrollo optimiza las dependencias Argon2 y Blake2 a nivel 3. Las cinco operaciones secuenciales con dos pasadas promediaron 400.3851816 ms, por debajo de la banda de 500 a 1000 ms. Con tres pasadas, el promedio de cinco operaciones fue 659.7596048 ms, dentro de esa banda. Se mantuvieron 256 MiB de memoria, una vía, sal aleatoria de 16 bytes y salida de 32 bytes. Esa base de dos pasadas por debajo de la banda motivó fijar tres pasadas para los hashes nuevos, según \rutarepo{docs/adr/0044-reference-password-hashing-cost.md}.

La comprobación focal aprobó siete casos, incluida la verificación de un PHC histórico de dos pasadas con contraseña correcta e incorrecta. Seis pruebas CLI aprobaron y la demostración completa de comandos terminó satisfactoriamente; su calibración secundaria dio 563.2 ms sobre cinco hashes. Las dos observaciones conservan su contexto. Las identidades del binario, entorno, comandos y resultados se registran en \rutarepo{docs/verification-report.md}.

El ajuste afecta a hashes nuevos; los anteriores conservan sus parámetros y no se reescriben. La banda no mide un inicio de sesión completo ni acredita rendimiento bajo carga. La memoria por operación no aumenta, aunque el trabajo de CPU y la ocupación del trabajador son mayores. Estas comprobaciones focales no constituyen la regresión global; la integración del ajuste en \texttt{main} se registra en \rutarepo{docs/product-completion.md}.
